\IfFileExists{stacks-project.cls}{%
\documentclass{stacks-project}
}{%
\documentclass{amsart}
}

% For dealing with references we use the comment environment
\usepackage{verbatim}
\newenvironment{reference}{\comment}{\endcomment}
%\newenvironment{reference}{}{}
\newenvironment{slogan}{\comment}{\endcomment}
\newenvironment{history}{\comment}{\endcomment}

% For commutative diagrams we use Xy-pic
\usepackage[all]{xy}

% We use 2cell for 2-commutative diagrams.
\xyoption{2cell}
\UseAllTwocells

% We use multicol for the list of chapters between chapters
\usepackage{multicol}

% This is generally recommended for better output
\usepackage{lmodern}
\usepackage[T1]{fontenc}

% For cross-file-references
\usepackage{xr-hyper}

% Package for hypertext links:
\usepackage{hyperref}

% For any local file, say "hello.tex" you want to link to please
% use \externaldocument[hello-]{hello}
\externaldocument[introduction-]{introduction}
\externaldocument[conventions-]{conventions}
\externaldocument[sets-]{sets}
\externaldocument[categories-]{categories}
\externaldocument[topology-]{topology}
\externaldocument[sheaves-]{sheaves}
\externaldocument[sites-]{sites}
\externaldocument[stacks-]{stacks}
\externaldocument[fields-]{fields}
\externaldocument[algebra-]{algebra}
\externaldocument[brauer-]{brauer}
\externaldocument[homology-]{homology}
\externaldocument[derived-]{derived}
\externaldocument[simplicial-]{simplicial}
\externaldocument[more-algebra-]{more-algebra}
\externaldocument[smoothing-]{smoothing}
\externaldocument[modules-]{modules}
\externaldocument[sites-modules-]{sites-modules}
\externaldocument[injectives-]{injectives}
\externaldocument[cohomology-]{cohomology}
\externaldocument[sites-cohomology-]{sites-cohomology}
\externaldocument[dga-]{dga}
\externaldocument[dpa-]{dpa}
\externaldocument[sdga-]{sdga}
\externaldocument[hypercovering-]{hypercovering}
\externaldocument[schemes-]{schemes}
\externaldocument[constructions-]{constructions}
\externaldocument[properties-]{properties}
\externaldocument[morphisms-]{morphisms}
\externaldocument[coherent-]{coherent}
\externaldocument[divisors-]{divisors}
\externaldocument[limits-]{limits}
\externaldocument[varieties-]{varieties}
\externaldocument[topologies-]{topologies}
\externaldocument[descent-]{descent}
\externaldocument[perfect-]{perfect}
\externaldocument[more-morphisms-]{more-morphisms}
\externaldocument[flat-]{flat}
\externaldocument[groupoids-]{groupoids}
\externaldocument[more-groupoids-]{more-groupoids}
\externaldocument[etale-]{etale}
\externaldocument[chow-]{chow}
\externaldocument[intersection-]{intersection}
\externaldocument[pic-]{pic}
\externaldocument[weil-]{weil}
\externaldocument[adequate-]{adequate}
\externaldocument[dualizing-]{dualizing}
\externaldocument[duality-]{duality}
\externaldocument[discriminant-]{discriminant}
\externaldocument[derham-]{derham}
\externaldocument[local-cohomology-]{local-cohomology}
\externaldocument[algebraization-]{algebraization}
\externaldocument[curves-]{curves}
\externaldocument[resolve-]{resolve}
\externaldocument[models-]{models}
\externaldocument[functors-]{functors}
\externaldocument[equiv-]{equiv}
\externaldocument[pione-]{pione}
\externaldocument[etale-cohomology-]{etale-cohomology}
\externaldocument[proetale-]{proetale}
\externaldocument[relative-cycles-]{relative-cycles}
\externaldocument[more-etale-]{more-etale}
\externaldocument[trace-]{trace}
\externaldocument[crystalline-]{crystalline}
\externaldocument[spaces-]{spaces}
\externaldocument[spaces-properties-]{spaces-properties}
\externaldocument[spaces-morphisms-]{spaces-morphisms}
\externaldocument[decent-spaces-]{decent-spaces}
\externaldocument[spaces-cohomology-]{spaces-cohomology}
\externaldocument[spaces-limits-]{spaces-limits}
\externaldocument[spaces-divisors-]{spaces-divisors}
\externaldocument[spaces-over-fields-]{spaces-over-fields}
\externaldocument[spaces-topologies-]{spaces-topologies}
\externaldocument[spaces-descent-]{spaces-descent}
\externaldocument[spaces-perfect-]{spaces-perfect}
\externaldocument[spaces-more-morphisms-]{spaces-more-morphisms}
\externaldocument[spaces-flat-]{spaces-flat}
\externaldocument[spaces-groupoids-]{spaces-groupoids}
\externaldocument[spaces-more-groupoids-]{spaces-more-groupoids}
\externaldocument[bootstrap-]{bootstrap}
\externaldocument[spaces-pushouts-]{spaces-pushouts}
\externaldocument[spaces-chow-]{spaces-chow}
\externaldocument[groupoids-quotients-]{groupoids-quotients}
\externaldocument[spaces-more-cohomology-]{spaces-more-cohomology}
\externaldocument[spaces-simplicial-]{spaces-simplicial}
\externaldocument[spaces-duality-]{spaces-duality}
\externaldocument[formal-spaces-]{formal-spaces}
\externaldocument[restricted-]{restricted}
\externaldocument[spaces-resolve-]{spaces-resolve}
\externaldocument[formal-defos-]{formal-defos}
\externaldocument[defos-]{defos}
\externaldocument[cotangent-]{cotangent}
\externaldocument[examples-defos-]{examples-defos}
\externaldocument[algebraic-]{algebraic}
\externaldocument[examples-stacks-]{examples-stacks}
\externaldocument[stacks-sheaves-]{stacks-sheaves}
\externaldocument[criteria-]{criteria}
\externaldocument[artin-]{artin}
\externaldocument[quot-]{quot}
\externaldocument[stacks-properties-]{stacks-properties}
\externaldocument[stacks-morphisms-]{stacks-morphisms}
\externaldocument[stacks-limits-]{stacks-limits}
\externaldocument[stacks-cohomology-]{stacks-cohomology}
\externaldocument[stacks-perfect-]{stacks-perfect}
\externaldocument[stacks-introduction-]{stacks-introduction}
\externaldocument[stacks-more-morphisms-]{stacks-more-morphisms}
\externaldocument[stacks-geometry-]{stacks-geometry}
\externaldocument[moduli-]{moduli}
\externaldocument[moduli-curves-]{moduli-curves}
\externaldocument[examples-]{examples}
\externaldocument[exercises-]{exercises}
\externaldocument[guide-]{guide}
\externaldocument[desirables-]{desirables}
\externaldocument[coding-]{coding}
\externaldocument[obsolete-]{obsolete}
\externaldocument[fdl-]{fdl}
\externaldocument[index-]{index}

% Theorem environments.
%
\theoremstyle{plain}
\newtheorem{theorem}[subsection]{Theorem}
\newtheorem{proposition}[subsection]{Proposition}
\newtheorem{lemma}[subsection]{Lemma}

\theoremstyle{definition}
\newtheorem{definition}[subsection]{Definition}
\newtheorem{example}[subsection]{Example}
\newtheorem{exercise}[subsection]{Exercise}
\newtheorem{situation}[subsection]{Situation}

\theoremstyle{remark}
\newtheorem{remark}[subsection]{Remark}
\newtheorem{remarks}[subsection]{Remarks}

\numberwithin{equation}{subsection}

% Macros
%
\def\lim{\mathop{\mathrm{lim}}\nolimits}
\def\colim{\mathop{\mathrm{colim}}\nolimits}
\def\Spec{\mathop{\mathrm{Spec}}}
\def\Hom{\mathop{\mathrm{Hom}}\nolimits}
\def\Ext{\mathop{\mathrm{Ext}}\nolimits}
\def\SheafHom{\mathop{\mathcal{H}\!\mathit{om}}\nolimits}
\def\SheafExt{\mathop{\mathcal{E}\!\mathit{xt}}\nolimits}
\def\Sch{\mathit{Sch}}
\def\Mor{\mathop{\mathrm{Mor}}\nolimits}
\def\Ob{\mathop{\mathrm{Ob}}\nolimits}
\def\Sh{\mathop{\mathit{Sh}}\nolimits}
\def\NL{\mathop{N\!L}\nolimits}
\def\CH{\mathop{\mathrm{CH}}\nolimits}
\def\proetale{{pro\text{-}\acute{e}tale}}
\def\etale{{\acute{e}tale}}
\def\QCoh{\mathit{QCoh}}
\def\Ker{\mathop{\mathrm{Ker}}}
\def\Im{\mathop{\mathrm{Im}}}
\def\Coker{\mathop{\mathrm{Coker}}}
\def\Coim{\mathop{\mathrm{Coim}}}

% Boxtimes
%
\DeclareMathSymbol{\boxtimes}{\mathbin}{AMSa}{"02}

%
% Macros for moduli stacks/spaces
%
\def\QCohstack{\mathcal{QC}\!\mathit{oh}}
\def\Cohstack{\mathcal{C}\!\mathit{oh}}
\def\Spacesstack{\mathcal{S}\!\mathit{paces}}
\def\Quotfunctor{\mathrm{Quot}}
\def\Hilbfunctor{\mathrm{Hilb}}
\def\Curvesstack{\mathcal{C}\!\mathit{urves}}
\def\Polarizedstack{\mathcal{P}\!\mathit{olarized}}
\def\Complexesstack{\mathcal{C}\!\mathit{omplexes}}
% \Pic is the operator that assigns to X its picard group, usage \Pic(X)
% \Picardstack_{X/B} denotes the Picard stack of X over B
% \Picardfunctor_{X/B} denotes the Picard functor of X over B
\def\Pic{\mathop{\mathrm{Pic}}\nolimits}
\def\Picardstack{\mathcal{P}\!\mathit{ic}}
\def\Picardfunctor{\mathrm{Pic}}
\def\Deformationcategory{\mathcal{D}\!\mathit{ef}}
\usepackage{cleveref}
\crefname{originalequation}{original equation}{original equations}
\Crefname{originalequation}{Original equation}{Original equations}
\makeatletter
\newlabel{equation-C}{{0AL3}{source}{}{sourcepointer-1}{}}
\newlabel{equation-C@cref}{{[originalequation][0][]0AL3}{[1][0][]source}{}{}{}}
\makeatother
\crefname{originalequation}{original equation}{original equations}
\Crefname{originalequation}{Original equation}{Original equations}
\makeatletter
\newlabel{equation-C-prime}{{0AL4}{source}{}{sourcepointer-2}{}}
\newlabel{equation-C-prime@cref}{{[originalequation][0][]0AL4}{[1][0][]source}{}{}{}}
\makeatother
\crefname{originalreference}{original reference}{original references}
\Crefname{originalreference}{Original reference}{Original references}
\makeatletter
\newlabel{etale-theorem-remarkable-equivalence}{{039R}{source}{}{sourcepointer-3}{}}
\newlabel{etale-theorem-remarkable-equivalence@cref}{{[originalreference][0][]039R}{[1][0][]source}{}{}{}}
\makeatother
\crefname{originallemma}{original lemma}{original lemmas}
\Crefname{originallemma}{Original lemma}{Original lemmas}
\makeatletter
\newlabel{lemma-NL-is-completion}{{0GAF}{source}{}{sourcepointer-4}{}}
\newlabel{lemma-NL-is-completion@cref}{{[originallemma][0][]0GAF}{[1][0][]source}{}{}{}}
\makeatother
\crefname{originallemma}{original lemma}{original lemmas}
\Crefname{originallemma}{Original lemma}{Original lemmas}
\makeatletter
\newlabel{lemma-approximate-principal}{{0ALS}{source}{}{sourcepointer-5}{}}
\newlabel{lemma-approximate-principal@cref}{{[originallemma][0][]0ALS}{[1][0][]source}{}{}{}}
\makeatother
\crefname{originallemma}{original lemma}{original lemmas}
\Crefname{originallemma}{Original lemma}{Original lemmas}
\makeatletter
\newlabel{lemma-base-change-rig-etale-homomorphism}{{0GB2}{source}{}{sourcepointer-6}{}}
\newlabel{lemma-base-change-rig-etale-homomorphism@cref}{{[originallemma][0][]0GB2}{[1][0][]source}{}{}{}}
\makeatother
\crefname{originallemma}{original lemma}{original lemmas}
\Crefname{originallemma}{Original lemma}{Original lemmas}
\makeatletter
\newlabel{lemma-equivalent-with-artin}{{0AJU}{source}{}{sourcepointer-7}{}}
\newlabel{lemma-equivalent-with-artin@cref}{{[originallemma][0][]0AJU}{[1][0][]source}{}{}{}}
\makeatother
\crefname{originallemma}{original lemma}{original lemmas}
\Crefname{originallemma}{Original lemma}{Original lemmas}
\makeatletter
\newlabel{lemma-fully-faithful-etale-over-complement}{{0AKJ}{source}{}{sourcepointer-8}{}}
\newlabel{lemma-fully-faithful-etale-over-complement@cref}{{[originallemma][0][]0AKJ}{[1][0][]source}{}{}{}}
\makeatother
\crefname{originallemma}{original lemma}{original lemmas}
\Crefname{originallemma}{Original lemma}{Original lemmas}
\makeatletter
\newlabel{lemma-lift-approximation}{{0ALT}{source}{}{sourcepointer-9}{}}
\newlabel{lemma-lift-approximation@cref}{{[originallemma][0][]0ALT}{[1][0][]source}{}{}{}}
\makeatother
\crefname{originallemma}{original lemma}{original lemmas}
\Crefname{originallemma}{Original lemma}{Original lemmas}
\makeatletter
\newlabel{lemma-rig-etale}{{0ALQ}{source}{}{sourcepointer-10}{}}
\newlabel{lemma-rig-etale@cref}{{[originallemma][0][]0ALQ}{[1][0][]source}{}{}{}}
\makeatother
\crefname{originallemma}{original lemma}{original lemmas}
\Crefname{originallemma}{Original lemma}{Original lemmas}
\makeatletter
\newlabel{lemma-topologically-finite-type}{{0AJP}{source}{}{sourcepointer-11}{}}
\newlabel{lemma-topologically-finite-type@cref}{{[originallemma][0][]0AJP}{[1][0][]source}{}{}{}}
\makeatother
\crefname{originalreference}{original reference}{original references}
\Crefname{originalreference}{Original reference}{Original references}
\makeatletter
\newlabel{more-morphisms-lemma-check-smoothness-on-infinitesimal-nbhds}{{0A43}{source}{}{sourcepointer-12}{}}
\newlabel{more-morphisms-lemma-check-smoothness-on-infinitesimal-nbhds@cref}{{[originalreference][0][]0A43}{[1][0][]source}{}{}{}}
\makeatother
\crefname{originalreference}{original reference}{original references}
\Crefname{originalreference}{Original reference}{Original references}
\makeatletter
\newlabel{morphisms-lemma-etale-permanence}{{02GW}{source}{}{sourcepointer-13}{}}
\newlabel{morphisms-lemma-etale-permanence@cref}{{[originalreference][0][]02GW}{[1][0][]source}{}{}{}}
\makeatother
\crefname{originalremark}{original remark}{original remarks}
\Crefname{originalremark}{Original remark}{Original remarks}
\makeatletter
\newlabel{remark-take-bar}{{0AL6}{source}{}{sourcepointer-14}{}}
\newlabel{remark-take-bar@cref}{{[originalremark][0][]0AL6}{[1][0][]source}{}{}{}}
\makeatother
\title{Rig-etale algebras in the category C-prime are completed finite type algebras}
\author{The Stacks Project Authors}
\date{stacks-project:restricted.tex:local-snapshot-20261012; GNU FDL 1.2 or later}
\begin{document}
\maketitle
\noindent\textit{Editorial note.} Counted span is the author's own proof environment, source lines 2316-2403, following the statement at 2307-2314.\par\medskip
\noindent\textit{Editorial note.} The single (details omitted) at line 2368 is the standard completion base change along an etale (flat) map, a one-step standard-fact instantiation covered by the card's obligations.\par\medskip
\noindent\textit{Editorial note.} The induction is on the number of generators of I; the reduction to the principal case is the organizing device of the whole proof.\par\medskip
\noindent\textit{Editorial note.} Steps the author leaves to the reader are flagged here for the reviewer: source line(s) 2368 carry an "omitted"/"follows easily"/"immediate from the definitions" phrase. Those sentences are retained verbatim; the displayed argument around them stays inside the counted span.\par\medskip
\noindent\textit{Allowed standard prerequisites (editorial).}\begin{itemize}
\item Principal ideal case. lemma-approximate-principal, cited at proof line 2320. Supplies the t = 1 base case and the final application to B' over (A, (a\_t)).
\item Base change and etale-over-complement control. lemma-base-change-rig-etale-homomorphism at proof lines 2328-2329 and lemma-rig-etale at lines 2336-2338. Supplies that B\_m = B/(a\_t\textasciicircum{}m) stays rig-etale over the reduced pair and that C\_m may be chosen etale over the complement of V(I-bar\_m).
\item Etale lifting machinery. lemma-fully-faithful-etale-over-complement (2350-2354), morphisms-lemma-etale-permanence (2357), more-morphisms-lemma-check-smoothness-on-infinitesimal-nbhds (2359-2360), etale-theorem-remarkable-equivalence (2362-2366), lemma-lift-approximation (2369-2378). Supplies the compatibilization step C\_m/(a\_t\textasciicircum{}\{m-1\}) -> C\_\{m-1\} and the lift C''\_m of C\_\{m-1\} to an algebra over A-bar\_m.
\item NL comparison and rig-etale transfer. lemma-topologically-finite-type (2383-2385), lemma-NL-is-completion (2392-2395), lemma-equivalent-with-artin (2396-2400). Supplies that B' = lim C\_m is topologically finite type, that the NL cohomologies agree, and the uniform annihilation transferring rig-etale-ness from B to B'.
\end{itemize}\medskip
\ifdefined\hypertarget\hypertarget{sourcepointer-1}{}\fi
\noindent\textit{Source reference.} Original equation 0AL3: sources/stacks-restricted/pinned-original-source.tex, source lines 312--312. Navigation pointer only: the counted proof cites this statement by \textbackslash{}ref as a standing input; the deduction the counted proof has to make is carried out inside the counted span, and no part of that deduction is delegated to this pointer.\par\medskip
\ifdefined\hypertarget\hypertarget{sourcepointer-2}{}\fi
\noindent\textit{Source reference.} Original equation 0AL4: sources/stacks-restricted/pinned-original-source.tex, source lines 326--326. Navigation pointer only: the counted proof cites this statement by \textbackslash{}ref as a standing input; the deduction the counted proof has to make is carried out inside the counted span, and no part of that deduction is delegated to this pointer.\par\medskip
\ifdefined\hypertarget\hypertarget{sourcepointer-3}{}\fi
\noindent\textit{Source reference.} Original reference 039R: focused-seeds/geometry/S9-restricted-approximate/ref-tags.txt, source lines 5--5. Navigation pointer only: the counted proof cites this statement by \textbackslash{}ref as a standing input; the deduction the counted proof has to make is carried out inside the counted span, and no part of that deduction is delegated to this pointer.\par\medskip
\ifdefined\hypertarget\hypertarget{sourcepointer-4}{}\fi
\noindent\textit{Source reference.} Original lemma 0GAF: sources/stacks-restricted/pinned-original-source.tex, source lines 560--560. Navigation pointer only: the counted proof cites this statement by \textbackslash{}ref as a standing input; the deduction the counted proof has to make is carried out inside the counted span, and no part of that deduction is delegated to this pointer.\par\medskip
\ifdefined\hypertarget\hypertarget{sourcepointer-5}{}\fi
\noindent\textit{Source reference.} Original lemma 0ALS: sources/stacks-restricted/pinned-original-source.tex, source lines 2281--2281. Navigation pointer only: the counted proof cites this statement by \textbackslash{}ref as a standing input; the deduction the counted proof has to make is carried out inside the counted span, and no part of that deduction is delegated to this pointer.\par\medskip
\ifdefined\hypertarget\hypertarget{sourcepointer-6}{}\fi
\noindent\textit{Source reference.} Original lemma 0GB2: sources/stacks-restricted/pinned-original-source.tex, source lines 2048--2048. Navigation pointer only: the counted proof cites this statement by \textbackslash{}ref as a standing input; the deduction the counted proof has to make is carried out inside the counted span, and no part of that deduction is delegated to this pointer.\par\medskip
\ifdefined\hypertarget\hypertarget{sourcepointer-7}{}\fi
\noindent\textit{Source reference.} Original lemma 0AJU: sources/stacks-restricted/pinned-original-source.tex, source lines 1810--1810. Navigation pointer only: the counted proof cites this statement by \textbackslash{}ref as a standing input; the deduction the counted proof has to make is carried out inside the counted span, and no part of that deduction is delegated to this pointer.\par\medskip
\ifdefined\hypertarget\hypertarget{sourcepointer-8}{}\fi
\noindent\textit{Source reference.} Original lemma 0AKJ: sources/stacks-restricted/pinned-original-source.tex, source lines 2065--2065. Navigation pointer only: the counted proof cites this statement by \textbackslash{}ref as a standing input; the deduction the counted proof has to make is carried out inside the counted span, and no part of that deduction is delegated to this pointer.\par\medskip
\ifdefined\hypertarget\hypertarget{sourcepointer-9}{}\fi
\noindent\textit{Source reference.} Original lemma 0ALT: sources/stacks-restricted/pinned-original-source.tex, source lines 2151--2151. Navigation pointer only: the counted proof cites this statement by \textbackslash{}ref as a standing input; the deduction the counted proof has to make is carried out inside the counted span, and no part of that deduction is delegated to this pointer.\par\medskip
\ifdefined\hypertarget\hypertarget{sourcepointer-10}{}\fi
\noindent\textit{Source reference.} Original lemma 0ALQ: sources/stacks-restricted/pinned-original-source.tex, source lines 1981--1981. Navigation pointer only: the counted proof cites this statement by \textbackslash{}ref as a standing input; the deduction the counted proof has to make is carried out inside the counted span, and no part of that deduction is delegated to this pointer.\par\medskip
\ifdefined\hypertarget\hypertarget{sourcepointer-11}{}\fi
\noindent\textit{Source reference.} Original lemma 0AJP: sources/stacks-restricted/pinned-original-source.tex, source lines 346--346. Navigation pointer only: the counted proof cites this statement by \textbackslash{}ref as a standing input; the deduction the counted proof has to make is carried out inside the counted span, and no part of that deduction is delegated to this pointer.\par\medskip
\ifdefined\hypertarget\hypertarget{sourcepointer-12}{}\fi
\noindent\textit{Source reference.} Original reference 0A43: focused-seeds/geometry/S9-restricted-approximate/ref-tags.txt, source lines 6--6. Navigation pointer only: the counted proof cites this statement by \textbackslash{}ref as a standing input; the deduction the counted proof has to make is carried out inside the counted span, and no part of that deduction is delegated to this pointer.\par\medskip
\ifdefined\hypertarget\hypertarget{sourcepointer-13}{}\fi
\noindent\textit{Source reference.} Original reference 02GW: focused-seeds/geometry/S9-restricted-approximate/ref-tags.txt, source lines 7--7. Navigation pointer only: the counted proof cites this statement by \textbackslash{}ref as a standing input; the deduction the counted proof has to make is carried out inside the counted span, and no part of that deduction is delegated to this pointer.\par\medskip
\ifdefined\hypertarget\hypertarget{sourcepointer-14}{}\fi
\noindent\textit{Source reference.} Original remark 0AL6: sources/stacks-restricted/pinned-original-source.tex, source lines 454--454. Navigation pointer only: the counted proof cites this statement by \textbackslash{}ref as a standing input; the deduction the counted proof has to make is carried out inside the counted span, and no part of that deduction is delegated to this pointer.\par\medskip
\begin{lemma}
\label{lemma-approximate}
Let $A$ be a Noetherian ring. Let $I \subset A$ be an ideal.
Let $B$ be an object of (\ref{equation-C-prime}) which is
rig-\'etale over $(A, I)$.
Then there exists a finite type $A$-algebra $C$ and an
isomorphism $B \cong C^\wedge$.
\end{lemma}

\begin{proof}
We prove this lemma by induction on the number of generators of $I$.
Say $I = (a_1, \ldots, a_t)$. If $t = 0$, then $I = 0$ and there
is nothing to prove. If $t = 1$, then the lemma follows from
Lemma \ref{lemma-approximate-principal}. Assume $t > 1$.

\medskip\noindent
For any $m \geq 1$ set $\bar A_m = A/(a_t^m)$. Consider the ideal
$\bar I_m = (\bar a_1, \ldots, \bar a_{t - 1})$ in $\bar A_m$.
Observe that $V(I \bar A_m) = V(\bar I_m)$.
Let $B_m = B/(a_t^m)$ be the base change of $B$ for the
map $(A, I) \to (\bar A_m, \bar I_m)$, see Remark \ref{remark-take-bar}.
By Lemma \ref{lemma-base-change-rig-etale-homomorphism}
we find that $B_m$ is rig-\'etale over $(\bar A_m, \bar I_m)$.

\medskip\noindent
By induction hypothesis (on $t$) we can find a finite type
$\bar A_m$-algebra $C_m$ and a map $C_m \to B_m$ which induces an
isomorphism $C_m^\wedge \cong B_m$
where the completion is with respect to $\bar I_m$.
By Lemma \ref{lemma-rig-etale} we may assume that
$\Spec(C_m) \to \Spec(\bar A_m)$ is \'etale
over $\Spec(\bar A_m) \setminus V(\bar I_m)$.

\medskip\noindent
We claim that we may choose $A_m \to C_m \to B_m$ as in the previous
paragraph such that moreover there are isomorphisms
$C_m/(a_t^{m - 1}) \to C_{m - 1}$ compatible with the given
$A$-algebra structure and the maps to $B_{m - 1} = B_m/(a_t^{m - 1})$.
Namely, first fix a choice of $A_1 \to C_1 \to B_1$.
Suppose we have found $C_{m - 1} \to C_{m - 2} \to \ldots \to C_1$
with the desired properties.
Note that $C_m/(a_t^{m - 1})$ is \'etale over
$\Spec(\bar A_{m - 1}) \setminus V(\bar I_{m - 1})$.
Hence by Lemma \ref{lemma-fully-faithful-etale-over-complement}
there exists an \'etale extension $C_{m - 1} \to C'_{m - 1}$
which induces an isomorphism modulo $\bar I_{m - 1}$ and an
$\bar A_{m - 1}$-algebra map $C_m/(a_t^{m - 1}) \to C'_{m - 1}$
inducing the isomorphism $B_m/(a_t^{m - 1}) \to B_{m - 1}$ on completions.
Note that $C_m/(a_t^{m - 1}) \to C'_{m - 1}$ is \'etale over the complement
of $V(\bar I_{m - 1})$ by
Morphisms, Lemma \ref{morphisms-lemma-etale-permanence}
and over $V(\bar I_{m - 1})$ induces an isomorphism on completions
hence is \'etale there too (for example by More on Morphisms, Lemma
\ref{more-morphisms-lemma-check-smoothness-on-infinitesimal-nbhds}).
Thus $C_m/(a_t^{m - 1}) \to C'_{m - 1}$ is \'etale. By the
topological invariance of \'etale morphisms
(\'Etale Morphisms, Theorem \ref{etale-theorem-remarkable-equivalence})
there exists an \'etale ring map $C_m \to C'_m$ such that
$C_m/(a_t^{m - 1}) \to C'_{m - 1}$ is isomorphic to
$C_m/(a_t^{m - 1}) \to C'_m/(a_t^{m - 1})$. Observe that the
$\bar I_m$-adic completion of $C'_m$ is equal to the $\bar I_m$-adic
completion of $C_m$, i.e., to $B_m$ (details omitted).
We apply Lemma \ref{lemma-lift-approximation} to the diagram
$$
\xymatrix{
 & C'_m \ar[r] & C'_m/(a_t^{m - 1}) \\
C''_m \ar@{..>}[ru] \ar@{..>}[rr] & & C_{m - 1} \ar[u] \\
 & \bar A_m \ar[r] \ar[uu] \ar@{..>}[lu] & \bar A_{m - 1} \ar[u]
}
$$
to see that there exists a ``lift'' of $C''_m$ of $C_{m - 1}$
to an algebra over $\bar A_m$ with all the desired properties.

\medskip\noindent
By construction $(C_m)$ is an object of the category
(\ref{equation-C}) for the principal ideal $(a_t)$.
Thus the inverse limit $B' = \lim C_m$ is an $(a_t)$-adically
complete $A$-algebra such that $B'/a_t B'$ is of finite type
over $A/(a_t)$, see Lemma \ref{lemma-topologically-finite-type}.
By construction we have $C_m = B'/(a_t^m)$, the $I$-adic
completion of $C_m$ is $B_m$, and the $I$-adic completion of
$B'$ is isomorphic to $B$.
For each $m$ the complex $\NL_{C_m/A_m}$ has finite cohomology
modules supported on $V(IC_m) \subset \Spec(C_m)$ by construction.
Hence these modules are $I$-adically complete (as these modules
are annihilated by a power of $I$). Since $\NL^\wedge_{B_m/A_m}$
is the $I$-adic completion of $\NL_{C_m/A_m}$, see
Lemma \ref{lemma-NL-is-completion}, it
follows that the map $\NL_{C_m/A_m} \to \NL_{B_m/A_m}$ induces
an isomorphism on cohomologies. Thus because $B$ is rig \'etale
over $A$ it follows that there is a fixed power of $I$ that annihilates
the cohomologies of $\NL_{C_m/A_m}$ for all $m$, see
Lemma \ref{lemma-equivalent-with-artin}. It follows that $B'$
is rig \'etale over $(A, (a_t))$ by the same lemma.
Hence finally, we may apply Lemma \ref{lemma-approximate-principal}
to $B'$ over $(A, (a_t))$ to finish the proof.
\end{proof}

\end{document}
